\documentclass{article}
\usepackage[UTF8]{ctex}
\usepackage{amsmath,amssymb,amsfonts}
\usepackage{enumitem}
\usepackage{geometry}
\geometry{a4paper, margin=1in}
% === 错题本排版优化（不参与题目标识） ===
\setlength{\parindent}{0pt}
\setlength{\parskip}{0.35em plus 0.10em minus 0.08em}
\setlength{\emergencystretch}{3em}
\setlist{itemsep=0.25em, topsep=0.40em, parsep=0pt}
\allowdisplaybreaks[2]
% === 错题本排版优化结束 ===
\begin{document}
\section*{中国矿业大学(北京)2024年高等代数真题}
\section*{22．中国矿业大学（北京）2024年高等代数试卷}

\subsection{题目1}
一．计算行列式．\\
（1）\(\left|\begin{array}{cccc}1 & 1 & 1 & 1 \\ x & y & z & w \\ x^{2} & y^{2} & z^{2} & w^{2} \\ x^{4} & y^{4} & z^{4} & w^{4}\end{array}\right|\) ．\\
（2）\(D_{n}=\left|\begin{array}{cccccc}2 a & 1 & 0 & \cdots & 0 & 0 \\ a^{2} & 2 a & 1 & \cdots & 0 & 0 \\ 0 & a^{2} & 2 a & \cdots & 0 & 0 \\ \vdots & \vdots & \vdots & & \vdots & \vdots \\ 0 & 0 & 0 & \cdots & 2 a & 1 \\ 0 & 0 & 0 & \cdots & a^{2} & 2 a\end{array}\right|_{n \times n}\)

\subsection{题目2}
二．设 \(a_{1}, a_{2}, \cdots, a_{n}\) 是 \(n\) 个两两互异的实数，\(f_{1}(x), f_{2}(x), \cdots, f_{n}(x)\) 是 \(n\) 个次数不大于 \(n-2\) 的实系数多项式，证明：

\[
\left|\begin{array}{cccc}
f_{1}\left(a_{1}\right) & f_{1}\left(a_{2}\right) & \cdots & f_{1}\left(a_{n}\right) \\
f_{2}\left(a_{1}\right) & f_{2}\left(a_{2}\right) & \cdots & f_{2}\left(a_{n}\right) \\
\vdots & \vdots & & \vdots \\
f_{n}\left(a_{1}\right) & f_{n}\left(a_{2}\right) & \cdots & f_{n}\left(a_{n}\right)
\end{array}\right|=0 \text {. }
\]
\end{document}
